\subsection{Operators}
\subsubsection{Selection}
Selection is performed using classic tournament selection, with
a tournament of size three. Three random individuals are selected
from the population and from those three the one with the highest
fitness is selected. Tournament selection helps prevent convergence
on local optima.
\subsubsection{Mutation}
Mutation is performed by selecting node probabilities to be mutated
with a $50\%$ chance. For any node selected, a random edge-probability
is chosen to be mutated by adding a random, normally distributed value,
with a standard deviation of $\sigma$, typically equal to $0.1$. Once
the value is added, all probabilities for the selected node are
re-normalized. Typically, mutation probability was set to $0.5$.
\subsubsection{Crossover} While crossover is not typically used
in evolution strategies, we felt it could help more quickly converge
on good solutions, especially combined with the $+$ operator.
Individuals are selected for crossover using tournament selection and
elitism. The selected individual and individual with the highest fitness
are used. 
Crossover is performed by iterating over the set of nodes and choosing,
with a $50\%$ chance, node probabilities to swap between the parents.
This crossover produces two new individuals. Typically, crossover
probability was set to $0.2$. 

\subsection{Experiment Comparison to Bazzan}
In \cite{bazzan2014evolutionary}, Bazzan et al. evolve traffic
assignments for paths discovered using the $k$-shortest paths
algorithm, with different values of $k$. Here we compare their results
with ours on the same problem, using our model. The significant
difference between our work and theirs is the use of $k$-shortest
paths to discover the fastest routes. We evolve traffic distributions
for each node. In their work they use origin-destination (OD) pairs from
two nodes, to two other nodes, creating four total OD pairs. For each OD
pair $k$ paths are discovered with $k$-shortest paths, and then assignments
are evolved. For consistency, we model the same graph by starting all
traffic at the origins they use, and assigning a safety of $1.0$ to each
destination, letting our model evolve traffic assignments through probabilities.
The results are compared below.
RESULTS
Our model has a longer average travel time than the results reported in
\cite{bazzan2014evolutionary}. We believe this is a result of starting
all the vehicles moving at the same time, in our model. However, we feel
our results are reasonable comparable, while offering a more robust model,
regarding dynamic events.


